\section{幂集。可数集}

\begin{enumerate}
\def\labelenumi{\arabic{enumi}.}
\tightlist
\item
  两个集合 E、F 若能被建立一一对应，则称它们等势。
\end{enumerate}

两个集合，若各自与第三个集合等势，则它们彼此等势。

若 E 与 F 等势，则 \(\mathfrak{P}(\mathbf{E})\) 并且 \(\mathfrak{P}(\mathbf{F})\) 等势。

若 E 与 F， \(E'\) 并且 \(F'\) , \(E''\) 并且 \(F''\) 分别等势，则 \(E \times E' \times E''\) 并且 \(F \times F' \times F''\) 等势。此命题可推广到任意多个集合的乘积。

\begin{enumerate}
\def\labelenumi{\arabic{enumi}.}
\setcounter{enumi}{1}
\tightlist
\item
  设 X 与 Y 是集合 E 的两个任意子集。关系“X 与 Y 等势”是 \(\mathfrak{P}(\mathbf{E})\) 上的一个等价关系。X 在此关系下所属的等价类称为 X 的幂（*），而这些类的集合（即 \(\mathfrak{P}(\mathbf{E})\) 按上述关系的商集）称为 E 的子集的幂的集合。
\end{enumerate}

若 E 与 F 是两个不同的集合，则 E 的子集 X 与 F 的子集 Y 之间的“X 与 Y 等势”关系可表述为：X 的幂与 Y 的幂等价。由此，我们在 E 的子集的幂的集合的一个子集与 F 的子集的幂的集合的一个子集之间建立了一一对应。

\begin{enumerate}
\def\labelenumi{\arabic{enumi}.}
\setcounter{enumi}{2}
\tightlist
\item
  设 E、F 为任意两个集合，它们可以相同也可以不同，并设 a（相应地 b）是 E（相应地 F）的子集的幂的集合中的一个元素。若存在一个从幂为 a 的子集 X ⊂ E 到幂为 b 的子集 Y ⊂ F 的一一映射，则称 a 小于 b，或 b 大于 a。若此外 a 与 b 不是等价的幂，则称 a 严格小于 b，或 b 严格大于 a。
\end{enumerate}

如果 \(\mathfrak{a}\) 并且 \(\mathfrak{b}\) 等价，则 \(\mathfrak{a}\) 既大于又小于 \(\mathfrak{b}\) 。反过来，有一个定理说，如果 \(\mathfrak{a}\) 既大于又小于 \(\mathfrak{b}\) ，那么 \(\mathfrak{a}\) 并且 \(\mathfrak{b}\) 是等价的。特别地，由此推出，一个集合 \(\mathbf{E}\) 的子集的势所成的集合按关系“ \(\mathfrak{a}\) 小于 \(\mathfrak{b}\) ”有序；当我们把这个集合称为有序集时，所指的总是由这个关系定义的序。

此外，利用佐恩引理（从而利用选择公理），可以证明集合 E 的子集的势所成的集合是良序的。

\begin{enumerate}
\def\labelenumi{\arabic{enumi}.}
\setcounter{enumi}{3}
\tightlist
\item
  一个集合 \(\mathbf{E}\) 的势严格小于集合 \(\mathfrak{P}(\mathbf{E})\) .
\end{enumerate}

如果 \(f\) 是把集合 \(\mathbf{E}\) 映入集合 \(\mathbf{F}\) 的一个映射，则任意子集的像 \(f(\mathbf{X})\) （此处 \(\mathbf{X}\) 是 \(\mathbf{E}\) 的子集）的势小于 \(\mathbf{X}\) 的势.

\begin{enumerate}
\def\labelenumi{\arabic{enumi}.}
\setcounter{enumi}{4}
\tightlist
\item
  令 \((\mathbf{X}_t)_{i \in I}\) 是集合 \(E\) 的子集族，使得 \(\mathbf{X}_t \cap \mathbf{X}_x = \emptyset\) （每当 \(i \neq x\) ）。设 \((\mathbf{Y}_t)_{i \in I}\) 是集合 \(F\) 的子集族，由同一集合 \(I\) 索引，并且使得 \(\mathbf{Y}_t\) 的势小于 \(\mathbf{X}_t\) 的势（对所有 \(i \in I\) ）。那么并集 \(\bigcup_{i \in J} \mathbf{Y}_t\) 的势小于 \(\bigcup_{i \in J} \mathbf{X}_t\) 的势（对所有子集 \(J\) （ \(I\) 的））.
\end{enumerate}

如果还有 \(Y_{t} \cap Y_{x} = \emptyset\) 每当 \(t \neq x\) ，并且如果 \(X_{t}\) 并且 \(Y_{t}\) 对所有……都是等势的 \(t \in I\) ，那么 \(\bigcup_{t \in J} X_{t}\) 与……等势 \(\bigcup_{t \in J} Y_{t}\) .

的势 \((\mathfrak{a}_{\iota})\) 特别地，若F与E是同一个集合，我们看到E的一族两两不相交子集的并集的势仅取决于这些子集的势；它被称为这些势的和（因此，仅当存在一族 \((\mathbf{X}_{\iota})\) 使得E的互不相交子集满足 \(X_{t}\) 具有幂 \(a_{t}\) 对于每个索引 t）。

如果 \((\mathbf{X}_i)_{i\in I}\) 并且 \((\mathbf{Y}_i)_{i\in I}\) 分别是E和F的子集族，由同一指标集I索引，使得 \(\mathbf{X}_i\) 并且 \(\mathbf{Y}_i\) 对所有……都是等势的 \(i\) ，则乘积 \(\prod_{i}\mathbf{X}_{i}\) 并且 \(\prod_{i}\mathbf{Y}_{i}\) 等势。

\begin{enumerate}
\def\labelenumi{\arabic{enumi}.}
\setcounter{enumi}{5}
\item
  自然整数集 N 可被视为某个无限集的有限子集的幂集。序关系“ \(x \leqslant y\) ”在 N 上正是对该幂集进行排序的关系，而两个自然整数之和是一个函数，与上述两个幂之和相同。
\item
  一个集合被称为可数的，如果它与自然整数集 N 的某个子集等势。因此，每个有限集都是可数的；若它有 n 个元素，则它与区间 \([0, n - 1]\) 集合N的。每个可数无限集都与N等势；特别地，N的每个无限子集都与N具有相同的基数。
\end{enumerate}

若E是无限集，则存在E的一个划分为可数无限集；特别地，每个无限集的基数都大于N的基数。

如果E是无限集，则集合 \(E \times E\) 并且 \(E \times N\) 都与E等势，且E的有限子集之集与E等势。特别地， \(N \times N\) 是一个可数无限集。

\begin{enumerate}
\def\labelenumi{\arabic{enumi}.}
\setcounter{enumi}{7}
\tightlist
\item
  集合E的元素序列，按定义是E的元素族，以集合N或N的子集为指标集。因此，指标集为N的序列写作 \((x_{n})_{n\in\mathbf{N}}\) ，或更简单地 \((x_{n})\) 当没有混淆可能时。若 n 表示一个一般整数， \(x_{n}\) 被称为数列的通项，或第n项。当n被替换为特定整数时，也使用后一种术语。数列的元素集合是可数的。
\end{enumerate}

一个序列被称为无限的或有限的，取决于其指标集是N的无限子集还是有限子集。有限序列的元素集合是有限的。

序列的每一个子族本身也是一个序列，称为给定序列的子序列。有限序列的每一个子序列都是有限序列。

一个元素族，其指标集为 \(N \times N\) ，或为 \(N \times N\) 称为双重序列。以 \(N \times N\) 被写成 \((x_{m,n})\) ，或更简单地 \((x_{mn})\) 为指标的双重序列，若无混淆风险，可简称为双重序列。对于具有两个以上指标的序列，情况类似。

两个序列 \((x_{n})\) , \((y_{n})\) 称为仅在其项的次序上有所不同，如果存在指标集的一个排列 \(f\) ，使得 \(y_{n} = x_{f(n)}\) （对所有 \(n\) ）.

对于任意元素族 \((x_{i})_{i\in I}\) （其指标集 \(I\) 为可数无穷集），我们可以如下关联一个无穷序列：存在一个双射 \(n\to f(n)\) （从 \(N\) 到 \(I\) 上的）；令 \(y_{n} = x_{f(n)}\) ，序列 \((y_{n})\) 称为把族 \((x_{i})\) 按 \(f\) 所定义的顺序加以排列而得到的。因此，对应于 \(N\) 到 \(I\) 上的两个不同双射的序列，仅在其项的次序上有所不同。

当 I 为有限时，以同样的方式操作，我们得到与族相关联的一个有限序列 \((x_{i})\) .

\begin{enumerate}
\def\labelenumi{\arabic{enumi}.}
\setcounter{enumi}{8}
\tightlist
\item
  一个族（family）的并、交和积 \((\mathbf{X}_t)_{t\in I}\) 的子集的 \(E\) 被称为可数的，如果 \(I\) 是一个可数集，有限的，如果 \(I\) 是有限的。
\end{enumerate}

如果 I 是可数的，并且若对所有 \(X_{t}\) ， \(t \in I\) 的幂小于给定的无限幂 a，则并集 \(\bigcup_{t} X_{t}\) 的幂小于 a。若此外

中至少有一个 \(X_{i}\) 具有幂 a，则 \(\bigcup_{i}X_{i}\) 具有幂 a。特别地，幂为 a 的集合的每个可数并集也具有幂 a；可数多个可数集合的并集是可数集合。

